%=============================================================================!
% 3.11  THE CURL TEST                                                         !
%=============================================================================!
\section{The curl test}
\label{sec:en-curl}

Given a force, how can we tell whether it is conservative? Checking the
work along every path is impossible. This section derives a test that
uses only derivatives of the force at each point, and explains what the
test measures: how much the force circulates around a small loop.

\subsection*{A property of every gradient}

\begin{toolbox}[Mixed partial derivatives]
For a function $U(x, y)$ whose second partial derivatives change smoothly
from place to place, the order of differentiation does not matter,
\begin{equation*}
   \frac{\partial}{\partial x}\Bigl(\frac{\partial U}{\partial y}\Bigr)
   = \frac{\partial}{\partial y}\Bigl(\frac{\partial U}{\partial x}\Bigr) .
\end{equation*}
For $U = x^2y^3$, for instance, $\partial U/\partial y = 3x^2y^2$, whose
derivative with respect to $x$ is $6xy^2$, and $\partial U/\partial x =
2xy^3$, whose derivative with respect to $y$ is also $6xy^2$. To see why
it holds in general, take the combination of the values of $U$ at the four
corners of a small rectangle,
\begin{equation*}
   D = U(x + \Delta x, y + \Delta y) - U(x + \Delta x, y)
   - U(x, y + \Delta y) + U(x, y) ,
\end{equation*}
and group its four terms in two ways. Grouped as
\begin{equation*}
   D = \bigl[U(x + \Delta x, y + \Delta y) - U(x + \Delta x, y)\bigr]
   - \bigl[U(x, y + \Delta y) - U(x, y)\bigr] ,
\end{equation*}
each bracket is a step north, one taken at $x + \Delta x$ and one at $x$.
Holding $x$ fixed, the fundamental theorem of calculus makes each bracket
the integral of $\partial U/\partial y$ along its step, so, writing
$\frac{\partial U}{\partial y}(a, b)$ for its value at the point $(a, b)$,
\begin{equation*}
   D = \int_y^{y + \Delta y}\Bigl[\frac{\partial U}{\partial y}(x + \Delta
   x, y') - \frac{\partial U}{\partial y}(x, y')\Bigr]\dd y'
\end{equation*}
exactly. By the same theorem, now holding $y'$ fixed, the square bracket
is the integral of $\partial(\partial U/\partial y)/\partial x$ from $x$ to
$x + \Delta x$. Because this second derivative changes smoothly, on a
small enough rectangle it stays as close as we like to its value at $(x,
y)$, so $D$ is $\Delta x\,\Delta y$ times $\partial(\partial U/\partial
y)/\partial x$ at $(x, y)$, with an error small compared with $\Delta
x\,\Delta y$. Grouped instead as
\begin{equation*}
   D = \bigl[U(x + \Delta x, y + \Delta y) - U(x, y + \Delta y)\bigr]
   - \bigl[U(x + \Delta x, y) - U(x, y)\bigr] ,
\end{equation*}
each bracket is a step east, and the same argument with $x$ and $y$
exchanged makes $D$ equal to $\Delta x\,\Delta y$ times $\partial(\partial
U/\partial x)/\partial y$ at $(x, y)$, again with an error small compared
with $\Delta x\,\Delta y$. Both expressions describe the same number $D$,
so dividing by $\Delta x\,\Delta y$ and letting the rectangle shrink, the
two mixed derivatives are equal.
\end{toolbox}

If $\vb{F} = -\grad U$ in the plane, then $F_x = -\partial U/\partial x$
and $F_y = -\partial U/\partial y$, so by the toolbox
\begin{equation*}
   \frac{\partial F_y}{\partial x}
   = -\frac{\partial}{\partial x}\Bigl(\frac{\partial U}{\partial y}\Bigr)
   = -\frac{\partial}{\partial y}\Bigl(\frac{\partial U}{\partial x}\Bigr)
   = \frac{\partial F_x}{\partial y} .
\end{equation*}
Every conservative force whose components have smoothly changing partial
derivatives therefore has
\begin{equation}
   \frac{\partial F_y}{\partial x} - \frac{\partial F_x}{\partial y} = 0
   \label{eq:en-curl-zero}
\end{equation}
at every point. The swirl $c\,(-y, x)$ of \cref{eq:en-swirl} fails: $F_y
= cx$ gives $\partial F_y/\partial x = c$, $F_x = -cy$ gives $\partial
F_x/\partial y = -c$, and the difference is $2c$. The swirl is not
conservative, as its work around a circle already showed.

\subsection*{Circulation around a small loop}

The combination in \cref{eq:en-curl-zero} has a meaning of its own, which
the next box finds, and which turns the necessary condition into a test.

\begin{toolbox}[The curl]
Take the small rectangle of \cref{fig:en-loop}, with corners $(x, y)$ and
$(x + \Delta x, y + \Delta y)$, and add up the work of a force $\vb{F}$
once around it anticlockwise, taking the force at the middle of each
side.
\begin{itemize}
\item Along the bottom, eastwards, the displacement is $(\Delta x, 0)$:
   work $F_x\,\Delta x$, with $F_x$ taken at height $y$.
\item Up the right side: $F_y\,\Delta y$, with $F_y$ taken at $x +
   \Delta x$.
\item Along the top, westwards, the displacement is $(-\Delta x, 0)$:
   work $-F_x\,\Delta x$, with $F_x$ taken at height $y + \Delta y$.
\item Down the left side: $-F_y\,\Delta y$, with $F_y$ taken at $x$.
\end{itemize}
Pairing the two sides that run east and west, and the two that run north
and south, the total is
\begin{align*}
   \oint\vb{F}\cdot\dd\vb{r}
   &\approx \bigl[F_y(x + \Delta x) - F_y(x)\bigr]\,\Delta y
   - \bigl[F_x(y + \Delta y) - F_x(y)\bigr]\,\Delta x\\
   &\approx \frac{\partial F_y}{\partial x}\,\Delta x\,\Delta y
   - \frac{\partial F_x}{\partial y}\,\Delta y\,\Delta x
   = \Bigl(\frac{\partial F_y}{\partial x}
   - \frac{\partial F_x}{\partial y}\Bigr)\Delta A ,
\end{align*}
where each square bracket is the change of one component over one short
step, and $\Delta A = \Delta x\,\Delta y$ is the area. The combination in
brackets is therefore the work around a small loop per unit area
enclosed. The work around a closed loop is called its circulation, so the
combination is the circulation per unit area; it is called the curl of
$\vb{F}$ in the plane. In three dimensions a small loop can lie in the
$xy$, the $yz$ or the $zx$ plane. Renaming the letters in turn, $x$ to
$y$, $y$ to $z$ and $z$ to $x$, carries a loop in the $xy$ plane,
traversed anticlockwise as seen from the positive end of the $z$ axis,
into one in the $yz$ plane, anticlockwise as seen from the positive end of
the $x$ axis, and carries $\partial F_y/\partial x - \partial F_x/\partial
y$ into $\partial F_z/\partial y - \partial F_y/\partial z$; renaming once
more gives the $zx$ plane and $\partial F_x/\partial z - \partial
F_z/\partial x$. The three combinations form the curl vector, written
$\grad\times\vb{F}$ and read `curl $\vb{F}$' (the $\times$ is the cross
product of Chapter~5, not needed here),
\begin{equation*}
   \grad\times\vb{F} = \Bigl(
   \frac{\partial F_z}{\partial y} - \frac{\partial F_y}{\partial z},\;
   \frac{\partial F_x}{\partial z} - \frac{\partial F_z}{\partial x},\;
   \frac{\partial F_y}{\partial x} - \frac{\partial F_x}{\partial y}
   \Bigr) ,
\end{equation*}
whose third component is the one found above, for a loop in the $xy$
plane.

A large loop made of the sides of many small rectangles, tiling the
region it encloses, has a circulation equal to the sum of the
circulations of the small rectangles. Each inner side belongs to two
neighbouring rectangles and is traversed once in each direction, so its
two contributions cancel, and only the outer boundary is left. Hence the
work around the loop is the sum of the curl times $\Delta A$ over the
rectangles. As the rectangles shrink, such a sum of values times small
areas approaches a limit, written $\iint\dots\dd A$ and called the
integral over the area, just as a sum of values times short steps
approaches the integral of \cref{sec:mo-integrating}; when the value is a
constant, the sum, and so the integral, is that constant times the whole
area. So
\begin{equation*}
   \oint\vb{F}\cdot\dd\vb{r} = \iint\Bigl(\frac{\partial F_y}{\partial x}
   - \frac{\partial F_x}{\partial y}\Bigr)\dd A ,
\end{equation*}
a result known as Green's theorem, for a loop traversed anticlockwise. It
needs $\vb{F}$ and its partial derivatives to change smoothly at every
point inside the loop, because every small rectangle of the tiling must
have a curl. A smooth curved loop is approached by staircases of ever
smaller rectangles. A staircase never becomes as short as the curve, which
would matter for friction; but for a force that depends only on position,
the work over each short stretch is nearly $\vb{F}\cdot\Delta\vb{r}$, and
the staircase and the stretch of curve beside it share the same
displacement $\Delta\vb{r}$, so the theorem holds for the curve too.
\end{toolbox}

\begin{figure}[htb]
\centering
\begin{tikzpicture}[line cap=round, line join=round, scale=1.1,
   side/.style={-{Stealth[length=5pt]}, line width=0.9pt, accent}]
   \draw[side] (0,0) -- (3,0) node[midway, below] {\small $F_x\,\Delta x$
      at $y$};
   \draw[side] (3,0) -- (3,2) node[midway, right] {\small $F_y\,\Delta y$
      at $x + \Delta x$};
   \draw[side] (3,2) -- (0,2) node[midway, above] {\small $-F_x\,\Delta x$
      at $y + \Delta y$};
   \draw[side] (0,2) -- (0,0) node[midway, left] {\small $-F_y\,\Delta y$
      at $x$};
   \fill (0,0) circle (1.3pt) node[below left] {\small $(x, y)$};
   \fill (3,2) circle (1.3pt) node[above right]
      {\small $(x + \Delta x, y + \Delta y)$};
\end{tikzpicture}
\caption{The work of a force once anticlockwise around a small rectangle,
side by side. The two sides that run east and west differ only in their
height, and the two that run north and south only in their east-west
position.}
\label{fig:en-loop}
\end{figure}

\subsection*{The test}

The curl test now reads in both directions. If a force is conservative,
its curl is zero everywhere, by \cref{eq:en-curl-zero}. Conversely, if
the curl is zero everywhere in a region with no holes, every loop in the
region encloses only points of the region, so Green's theorem makes the
work around it zero; within the region the work then does not depend on
the path, and the force is conservative there. In three dimensions
the same holds, because any loop can be spanned by a surface tiled with
small loops in the same way, a result called Stokes' theorem whose details
we do not need. The proviso about holes is real: \cref{ex:en-vortex} gives
a force whose curl is zero wherever it is defined, but which is not
defined at the origin, and whose work around the origin is not zero.

For the swirl, with curl $2c$ everywhere, Green's theorem gives the work
anticlockwise around any loop as $2c$ times the area it encloses. On the
circle of the line-integral toolbox that is $2c\times\pi R^2 = 2\pi
cR^2$, as computed there directly.

When the test is passed, $U$ follows from \cref{eq:en-potential-3d} along
whichever path is easiest. From the origin, go along the $x$ axis to $(x,
0)$, a leg on which only $F_x$ does work, and then straight north to $(x,
y)$, on which only $F_y$ does:
\begin{equation*}
   U(x, y) = -\int_0^x F_x(x', 0)\,\dd x' - \int_0^y F_y(x, y')\,\dd y' .
\end{equation*}

For many atoms the test has the same form, with every pair of the $3N$
coordinates in place of $x$ and $y$. Numbering the coordinates $q_1,
\dots, q_{3N}$, a force is conservative only if
\begin{equation*}
   \frac{\partial F_a}{\partial q_b} = \frac{\partial F_b}{\partial q_a}
   \qquad\text{for every pair $a$ and $b$,}
\end{equation*}
where $F_a$ is the force component along $q_a$. For a conservative
force these derivatives are minus the second derivatives of $U$, which
Chapter~4 collects into a table called the Hessian.

\notebook{03}{compute the curl of a force symbolically, and the work
around loops}

\begin{selfcheck}
Is the force $\vb{F} = (y, x)$ conservative? If so, find a potential
energy $U$ with $\vb{F} = -\grad U$.
\end{selfcheck}
